\chapter{Expanded Chapter Text Blocks for Final Integration}

The following text blocks can be moved into the main chapters during final editing if a longer body is required. They are written in thesis style and avoid references to appendices.

\section{L - Precision agriculture as a systems problem}

Precision agriculture should not be reduced to the use of sensors. It is a systems problem that requires the integration of agronomy, electronics, communication, software engineering, data science, and human decision-making. A sensor without a decision workflow only produces data. A machine learning model without reliable data only produces uncertain predictions. A dashboard without actuation remains a monitoring interface. A pump without safety logic can create operational risk. The value of a Digital Twin irrigation system emerges when these components are connected in a coherent cycle that supports reliable decisions.

\section{L - Why MQTT is appropriate for Agrio}

The publish/subscribe model of MQTT is appropriate for Agrio because it decouples the producers and consumers of data. The Raspberry Pi gateway does not need to know which services will consume the data; it only publishes to the sensors/data topic. The backend subscribes to that topic and processes messages. Similarly, the backend does not directly open a physical connection to the WeMos; it publishes commands to the command topic, and the WeMos subscribes. This architecture is flexible because additional subscribers, such as logging services or mobile notification services, can be added without modifying the sensor station.

\section{L - Why a relational database is useful}

A relational database is useful in Agrio because the system contains many linked entities: farms, plots, zones, sensors, readings, weather records, states, predictions, simulations, recommendations, commands, events, and feedback. Relational modeling makes these links explicit through primary keys and foreign keys. It also supports traceability, which is essential for validating decisions. For example, a recommendation can be linked to the zone state and prediction that produced it, and a command can be linked to the event that triggered it.

\section{L - Why Random Forest is a pragmatic first model}

Random Forest is a pragmatic first model for a final-year smart irrigation prototype because it works well with tabular data, handles nonlinear relations, and does not require the large sequential datasets needed by deep learning models. It is also more robust than a single decision tree because it averages several trees. However, Random Forest is not a substitute for agronomic validation. Its predictions depend on the quality and representativeness of the dataset. Therefore, in Agrio, Random Forest should be interpreted as a decision-support component rather than an absolute agronomic authority.

\section{L - Why human validation matters}

Human validation matters because irrigation decisions have physical consequences. A false positive decision can waste water or cause over-irrigation, while a false negative decision can stress plants. In a prototype, sensor calibration, model uncertainty, and communication reliability are still being evaluated. Keeping the farmer in the loop makes the system safer and more acceptable. It also generates a useful decision history that can later be used to improve the model.

\section{L - Why the prototype is still valuable without field-scale validation}

A controlled prototype cannot prove full agronomic performance, but it can prove system feasibility. It demonstrates that the data chain works, that the backend can process real payloads, that the database supports traceability, that the ML service can be integrated, that the Digital Twin can simulate scenarios, that the dashboard can display decisions, and that the pump can be controlled. These are necessary steps before any real field deployment.
